% =============================================================
%  第四单元  乘法口诀（一）
% =============================================================
\Chapter{第四关\quad 编顺口溜的人}
\noindent{\small\kaishu\color{tipc}（课本第四单元：乘法口诀（一），P43--P52）}

\begin{blunder}
豆角爸爸小学二年级，把“六七四十二”背成了“六七四十三”，
被罚抄五十遍。抄完豆角爸爸从此一见口诀就烦，一直烦到现在。
\textbf{所以这一关，豆角爸爸一个字都不让小豆角抄。}小豆角要是抄，说明豆角爸爸教错了。
\end{blunder}

\begin{gamble}{赌一把：对口令，输的人学猫叫}
\begin{itemize}
  \item 豆角爸爸说前半句，小豆角答后半句；说错就换小豆角出题。
  \item \textbf{豆角爸爸故意答错两次}（“四五……二十五！”），让小豆角来纠。
  \item 连错三次，学三声猫叫，声音要像。\textbf{豆角爸爸会故意输一次}，让小豆角听豆角爸爸叫。
\end{itemize}
\end{gamble}

\begin{bonus}
\begin{itemize}
  \item \textbf{把口诀唱出来}：用《两只老虎》的调子唱“一五得五、二五一十”。
        难听没关系，\textbf{越难听记得越牢}。允许小豆角改词、加节奏、配动作。
  \item \textbf{饭前抢答}：上桌前豆角爸爸喊一句前半句，谁先答出后半句\textbf{谁可以不吃青菜}。
        豆角爸爸故意慢半拍，让小豆角赢。
  \item \textbf{镜子卡片}：把小豆角总错的三五句写在卡片上贴冰箱上，
        拿酸奶背一遍、拿饮料背一遍、看课程表又背一遍。\textbf{只贴错的那几句，别贴整张表。}
\end{itemize}
\end{bonus}

\begin{unitinfo}
\textbf{课本 P43--P52}\quad \textbf{5 次课，每次 20 分钟}\quad
\textbf{建议第 11--13 周} \\[4pt]
\textbf{这一关到底在干嘛}：让小豆角\textbf{自己把口诀编出来}，不是背下来。
编过一遍，后面 2、3、4 甚至 6、7、8、9 的口诀小豆角都会自己编。\\
\textbf{逻辑是这么走的}：\ 摆出图 $\to$ 数出结果 $\to$ 压成一句顺口的话
$\to$ 发现相邻两句差一个固定数 $\to$ 忘了能推。\\
\textbf{一句话}：口诀不是天上掉下来的，是\textbf{从图里数出来的}。
跳过“数”这一步，45 句就是死记；走过这一步，45 句只有 9 句要记。
\end{unitinfo}

\subsection{开局先聊个事：口诀为什么长这样}

“三五十五”这四个字，其实是三件事压在一起：
\textbf{三个、五个、十五}。古人嫌“$5+5+5$”写起来麻烦，就编成了顺口溜。

豆角爸爸注意一个细节：为什么“一五得五”有个“得”，而“二五一十”没有？
因为“五”是一位数，念起来不顺，加个“得”垫一下；
“一十”是两位数，不加“得”念着也顺。\textbf{这是说话的顺口，不是规定。}

所以别让小豆角去背“什么时候加得、什么时候不加”，让小豆角\textbf{念一遍听顺不顺}就知道了。

\subsection{任务 1：把 5 的口诀编出来}

\begin{lesson}{1}{20}
\textbf{今天要拿下的}：九句 5 的口诀，一句一句从手上数出来。

\begin{materials}
小豆角自己的两只手、百数表（附录 B）、笔。
\end{materials}

\begin{enumerate}
  \item \textbf{数手}：一双手 5 个手指。一双 5 个，两双 10 个……
        一双一双加到 9 双。
  \item \textbf{填表}（让小豆角自己填，不会就 5 个 5 个加）：
        \begin{center}
        \small
        \begin{tabular}{@{}l*{9}{c}@{}}
        \toprule
        有几双 & 1 & 2 & 3 & 4 & 5 & 6 & 7 & 8 & 9 \\
        \midrule
        有几个 & 5 & 10 & 15 & 20 & 25 & 30 & 35 & 40 & 45 \\
        \bottomrule
        \end{tabular}
        \end{center}
  \item \textbf{一句一句编}（\textbf{别一次给九句}，一句一句来）：
        1 个 5 是 5 $\to$ “一五得五”；2 个 5 是 10 $\to$ “二五一十”；
        3 个 5 是 15 $\to$ “三五十五”；4 个 5 是 20 $\to$ “四五二十”……
  \item \textbf{让小豆角自己发现规律}（这一步别替小豆角总结）：
        \begin{itemize}
          \item 每句开头是“几个”，后面是结果。
          \item 一位数加“得”，两位数不加（因为念着顺）。
          \item \textbf{每句的最后一个字，不是“五”就是“十”。}
        \end{itemize}
  \item \textbf{抄成一张表}贴起来：
        \begin{center}
        \small
        \begin{tabular}{@{}ll@{\hspace{2.2cm}}ll@{}}
        \toprule
        一五得五 & $5\times1=5$ & 五六三十 & $5\times6=30$ \\
        二五一十 & $5\times2=10$ & 五七三十五 & $5\times7=35$ \\
        三五十五 & $5\times3=15$ & 五八四十 & $5\times8=40$ \\
        四五二十 & $5\times4=20$ & 五九四十五 & $5\times9=45$ \\
        五五二十五 & $5\times5=25$ & & \\
        \bottomrule
        \end{tabular}
        \end{center}
        \textit{口诀习惯小数在前，写“五六三十”，不写“六五三十”——念着顺而已。}
  \item \textbf{玩“挑错”}：豆角爸爸故意念错一句（“三五十二”），小豆角得立刻喊
        “不对，三五十五”。\textbf{小豆角纠正豆角爸爸时最来劲。}
\end{enumerate}

\begin{saybox}
“小豆角怎么知道三五就是十五？除了背，还能怎么知道？”（$5+5+5$）\\
“一五得五、二五一十，为什么第二句多了个‘一’？”（因为 10 是两个数字）\\
“五八四十，那五九多少？小豆角怎么想的？”（再加一个 5）\\
“这些口诀最后一个字有什么规律？”（不是“五”就是“十”）
\end{saybox}

\begin{passbox}
\begin{itemize}
  \item 9 句 5 的口诀，看着表小豆角能一句一句编出来。
  \item 小豆角能说出“口诀是从一个一个数出来的”。
\end{itemize}
\end{passbox}
\end{lesson}

\subsection{任务 2：忘了怎么办 + 编出 2 的口诀}

\begin{lesson}{2}{20}
\textbf{今天要拿下的}：\textbf{忘了一句能自己推出来}——这是本章最值钱的技能。

\begin{materials}
百数表、彩笔、两根小棒（当一双）。
\end{materials}

\begin{enumerate}
  \item \textbf{教推法}（比背重要十倍）：
        \begin{itemize}
          \item 忘了“五七三十五”？想前一句“五六三十”，$30+5=35$。
          \item 忘了“四五二十”？想后一句“五五二十五”，$25-5=20$。
        \end{itemize}
        \textbf{让小豆角故意“忘”一句，然后用前后推。}练三次。
        \textbf{会推的小豆角不怕忘}，不会推的小豆角一忘就慌。
  \item \textbf{百数表上圈 5 的倍数}：把 5、10、15……45 圈出来，
        看排成什么形状（都在最右一列或斜着下来）。
        \textbf{这就是“5 的倍数末位是 0 或 5”的来处}，不用背。
  \item \textbf{换材料编 2 的口诀}：两根小棒是一双。
        \begin{center}
        \small
        \begin{tabular}{@{}l*{9}{c}@{}}
        \toprule
        有几双 & 1 & 2 & 3 & 4 & 5 & 6 & 7 & 8 & 9 \\
        \midrule
        有几根 & 2 & 4 & 6 & 8 & 10 & 12 & 14 & 16 & 18 \\
        \bottomrule
        \end{tabular}
        \end{center}
        编：一二得二、二二得四、二三得六、二四得八、二五一十、
        二六十二、二七十四、二八十六、二九十八。
  \item \textbf{一个发现}：“二五一十”既能算 2 个 5，也能算 5 个 2。
        \textbf{一句话管两个算式}——这就是口诀省事的地方。
  \item \textbf{规律小结}：2 的口诀\textbf{得数全是双数}，相邻两句差 2；
        5 的口诀得数\textbf{末位是 0 或 5}，相邻差 5。
        让小豆角在百数表上圈 2 的倍数，自己看形状。
\end{enumerate}

\begin{saybox}
“五七三十五忘了，小豆角打算怎么想起来？”（加 5 或减 5）\\
“为什么 2 的口诀得数全是双数？”\\
“二五一十这句话，能算 2 个 5，还能算几个几？”（5 个 2）
\end{saybox}

\begin{passbox}
\begin{itemize}
  \item 5 和 2 的口诀，不看书能顺序说、也能跳着说。
  \item \textbf{故意说“豆角爸爸忘了”，小豆角能用加减推出来。}
\end{itemize}
\end{passbox}
\end{lesson}

\subsection{任务 3：让小豆角自己编 3 和 4 的口诀}

\begin{lesson}{3}{20}
\textbf{今天要拿下的}：这一次豆角爸爸不出声，让小豆角自己把 3、4 编出来。

\begin{materials}
小棒（摆三角形和正方形）、笔、纸。
\end{materials}

\begin{enumerate}
  \item \textbf{给情境，不给口诀}：“一辆三轮车 3 个轮子。2 辆呢？3 辆呢？”
        （课本 P47）
        \begin{center}
        \small
        \begin{tabular}{@{}l*{9}{c}@{}}
        \toprule
        有几辆 & 1 & 2 & 3 & 4 & 5 & 6 & 7 & 8 & 9 \\
        \midrule
        几个轮子 & 3 & 6 & 9 & 12 & 15 & 18 & 21 & 24 & 27 \\
        \bottomrule
        \end{tabular}
        \end{center}
  \item \textbf{让小豆角自己编}：豆角爸爸\textbf{全程不出声}，让小豆角照 5 的口诀的样子编。
        编错了也别急着纠，让小豆角自己念一遍听不顺。
        参考：一三得三、二三得六、三三得九、三四十二、三五十五、
        三六十八、三七二十一、三八二十四、三九二十七。
  \item \textbf{换 4}：一辆小汽车 4 个轮子。
        \begin{center}
        \small
        \begin{tabular}{@{}l*{9}{c}@{}}
        \toprule
        有几辆 & 1 & 2 & 3 & 4 & 5 & 6 & 7 & 8 & 9 \\
        \midrule
        几个轮子 & 4 & 8 & 12 & 16 & 20 & 24 & 28 & 32 & 36 \\
        \bottomrule
        \end{tabular}
        \end{center}
        口诀：一四得四、二四得八、三四十二、四四十六、四五二十、
        四六二十四、四七二十八、四八三十二、四九三十六。
  \item \textbf{找“一句管两个算式”}：
        \begin{center}
        \small
        \begin{tabular}{@{}lll@{}}
        \toprule
        \textbf{口诀} & \textbf{能算的算式} & \textbf{几个} \\
        \midrule
        三四十二 & $3\times4=12$、$4\times3=12$ & 两个 \\
        四五二十 & $4\times5=20$、$5\times4=20$ & 两个 \\
        四四十六 & $4\times4=16$ & \textbf{只有一个}（两个乘数一样）\\
        三三得九 & $3\times3=9$ & \textbf{只有一个} \\
        \bottomrule
        \end{tabular}
        \end{center}
  \item \textbf{教三招应对“忘了”}（写在纸上贴起来）：
        ① 想前一句加一个；② 想后一句减一个；③ 画图数一数。
\end{enumerate}

\begin{saybox}
“这句为什么叫‘三四十二’不叫‘四三十二’？”（小数在前，念着顺）\\
“三四十二除了算 $3\times4$，还能算哪道？”（$4\times3$）\\
“小豆角编的和书上一样吗？哪儿不一样？”
\end{saybox}

\begin{pitbox}
\textbf{两个常见问题}：
\begin{enumerate}[label=\arabic*、]
  \item \textbf{字写错}：写成“三四一十二”“三五五”。
        \textbf{解法}：一位数加“得”，两位数的“十二”别写成“一十二”。
        让小豆角自己念一遍，念着别扭就是写错了。
  \item \textbf{只会顺着背}。问“四五”要停三秒，问“五九”直接卡住。
        \textbf{解法}：抽查一律\textbf{跳着问}（三五、四八、二七），别按顺序。
\end{enumerate}
\end{pitbox}

\begin{passbox}
\begin{itemize}
  \item 3 和 4 的口诀，不看书能自己编出来。
  \item 小豆角知道“一句口诀能算两道乘法”（两个乘数不一样时）。
  \item 跳着问 10 句，对 8 句以上。
\end{itemize}
\end{passbox}
\end{lesson}

\subsection{任务 4：对口令，找规律}

\begin{lesson}{4}{20}
\textbf{今天要拿下的}：练到“不用想”，并且能从一串数里看出规律。

\begin{materials}
附录 A 的乘法卡（先用 1--5 那部分）。
\end{materials}

\begin{enumerate}
  \item \textbf{对口令}（课本 P44，最省事的复习法）：
        \begin{playbox}{对口令}
        豆角爸爸说前半句，小豆角说后半句：豆角爸爸说“四五”，小豆角说“二十”；豆角爸爸说“三七”，小豆角说“二十一”。\\
        然后\textbf{交换}：小豆角说前半句豆角爸爸说后半句——\textbf{豆角爸爸故意答错一两次}，
        让小豆角来纠豆角爸爸。小豆角会非常投入。\\
        进阶：只说数字，“3 和 4 相乘”$\to$ 小豆角答“十二”。
        \end{playbox}
  \item \textbf{找规律填数}（课本 P48）：
        \begin{enumerate}[label=（\arabic*）,itemsep=3pt]
          \item $10,\,15,\,20,$ \blank[1.4cm]，\blank[1.4cm]。（每次 $+5$ $\to$ 25、30）
          \item $18,\,16,\,14,$ \blank[1.4cm]，\blank[1.4cm]。（每次 $-2$ $\to$ 12、10）
          \item $20,\,24,\,28,$ \blank[1.4cm]，\blank[1.4cm]。（每次 $+4$ $\to$ 32、36）
          \item $27,\,24,\,21,$ \blank[1.4cm]，\blank[1.4cm]。（每次 $-3$ $\to$ 18、15）
        \end{enumerate}
        \textbf{每填一组，都要小豆角说出“每次加几／减几”}，再说出对应的口诀。
  \item \textbf{找朋友}：左边算式，右边口诀，打乱后连线。
        \begin{center}
        \small
        \begin{tabular}{@{}llll@{}}
        $8\times2$ & 二八十六 & $3\times9$ & 三九二十七 \\
        $6\times3$ & 三六十八 & $5\times4$ & 四五二十 \\
        \end{tabular}
        \end{center}
  \item \textbf{30 秒速答}：抽卡，30 秒内对 15 张算过。
  \item \textbf{错的单独挑出来}：哪几句总错，写卡片背面，睡前念三遍。
        \textbf{只念错的那几句，别整表重背}——整表重背是浪费时间。
\end{enumerate}

\begin{passbox}
\begin{itemize}
  \item 对口令反应在 2 秒以内。
  \item 四组找规律全对，能说出“每次加几／减几”。
  \item 30 秒对 15 张卡。
\end{itemize}
\end{passbox}
\end{lesson}

\subsection{任务 5：用口诀去解决事}

\begin{lesson}{5}{20}
\textbf{今天要拿下的}：遇到乘法问题，先想“用哪句口诀”，再写算式。

\begin{materials}
纸、笔。
\end{materials}

\begin{enumerate}
  \item \textbf{课本 P49 回家路上}：船上有 4 排小猫，每排 4 只。
        一共几只？$\to$ “四四十六” $\to 4\times4=16$（只）。
  \item \textbf{草地上的花}：2 组，每组 5 朵 $\to$ “二五一十” $\to 10$（朵）。
  \item \textbf{树上的鸟}：3 组，每组 5 只 $\to$ “三五十五” $\to 15$（只）。
  \item \textbf{反过来：给算式找事}。“$2\times6=12$ 能解决什么问题？”
        让小豆角想两个，\textbf{要求说得出的故事必须真的对得上算式}。
  \item \textbf{四步法}（贴在桌上）：
        \begin{center}
        \small
        \begin{tabularx}{\linewidth}{@{}p{1.0cm}p{1.8cm}X@{}}
        \toprule
        \textbf{步} & \textbf{干什么} & \textbf{例子：每层 3 米，7 层楼多高？} \\
        \midrule
        1 & 圈数字 & 3、7 \\
        2 & 说关系 & 每层 3 米，7 层就是 7 个 3 \\
        3 & 想口诀 & 三七二十一 \\
        4 & 写完整 & $3\times7=21$（米）。答：这楼高约 21 米。 \\
        \bottomrule
        \end{tabularx}
        \end{center}
  \item \textbf{自己做四道}（课本 P49--P50）：
        ① 一辆缆车坐 2 人，7 辆坐多少人？\\
        ② 5 个小丑，每人 3 个球，一共几个？\\
        ③ 每盒 4 个，6 盒一共几个？\\
        ④ 每排 5 人，4 排一共几人？
\end{enumerate}

\begin{saybox}
“这题用哪句口诀？小豆角怎么想到的？”\\
“算完了，单位是‘只’还是‘排’？”\\
“$3\times8$ 小豆角能编一个家里的事吗？”（比如一盒鸡蛋 3 排、每排 8 个……）
\end{saybox}

\begin{passbox}
\begin{itemize}
  \item 4 道应用题全对，\textbf{算式、单位、答语都齐}。
  \item 小豆角能说出用的是哪一句口诀。
\end{itemize}
\end{passbox}
\end{lesson}

\subsection{1--5 口诀总表（复印一张贴墙上）}

\begin{center}
\small
\begin{tabular}{@{}p{4.2cm}p{4.2cm}p{4.2cm}@{}}
\toprule
\textbf{5 的口诀} & \textbf{2 的口诀} & \textbf{3 的口诀} \\
\midrule
一五得五 & 一二得二 & 一三得三 \\
二五一十 & 二二得四 & 二三得六 \\
三五十五 & 二三得六 & 三三得九 \\
四五二十 & 二四得八 & 三四十二 \\
五五二十五 & 二五一十 & 三五十五 \\
五六三十 & 二六十二 & 三六十八 \\
五七三十五 & 二七十四 & 三七二十一 \\
五八四十 & 二八十六 & 三八二十四 \\
五九四十五 & 二九十八 & 三九二十七 \\
\midrule
\multicolumn{3}{@{}l}{\textbf{4 的口诀：}%
一四得四\quad 二四得八\quad 三四十二\quad 四四十六\quad 四五二十} \\
\multicolumn{3}{@{}l}{\hspace*{5.4em}%
四六二十四\quad 四七二十八\quad 四八三十二\quad 四九三十六} \\
\bottomrule
\end{tabular}
\end{center}

\subsection{翻车现场（对着看）}

\begin{center}
\small
\begin{tabularx}{\linewidth}{@{}p{3.8cm}p{4.0cm}X@{}}
\toprule
\textbf{小豆角写出来的样子} & \textbf{其实是想错了什么} & \textbf{豆角爸爸回一句什么} \\
\midrule
“三五五”“三四一十二” & 没听顺口 & 念一遍，念着别扭就是错了。 \\
\addlinespace[2pt]
只会顺着背，跳着问就卡 & 靠顺序记，没靠关系 & 抽查必须跳着问。 \\
\addlinespace[2pt]
忘一句就整段崩 & 不会用加减推 & “前一句加一个／后一句减一个。” \\
\addlinespace[2pt]
“三四十二”只算 $3\times4$ & 不知道一句管两道 & “倒过来也行，$4\times3$。” \\
\addlinespace[2pt]
单位写错（只／排） & 只盯数字不看问题 & 写完算式回头读一遍问题。 \\
\addlinespace[2pt]
算式对，口诀说不出 & 算和背脱节 & 每题都要报“用的哪句口诀”。 \\
\bottomrule
\end{tabularx}
\end{center}

\subsection{抽查：随机开枪（10 题）}

\begin{enumerate}[label=\arabic*.,itemsep=4pt]
  \item 三五 \blank[1.6cm] \qquad 2. 四六 \blank[1.6cm] \qquad 3. 二七 \blank[1.6cm]
  \item $5\times8=$ \blank[1.4cm] \qquad 5. $3\times9=$ \blank[1.4cm] \qquad
  6. $4\times7=$ \blank[1.4cm]
  \item $2\times6=$ \blank[1.4cm] \qquad 8. $5\times5=$ \blank[1.4cm]
  \item 每盒 4 个，7 盒一共 \blank[1.4cm] 个。
  \item 一辆车坐 5 人，6 辆坐 \blank[1.4cm] 人。
\end{enumerate}

\begin{ansbox}
1. 十五\quad 2. 二十四\quad 3. 十四\quad 4. 40\quad 5. 27\quad
6. 28\quad 7. 12\quad 8. 25\quad 9. 28\quad 10. 30 \\
\textbf{错 3 题以上}：别进第五单元，把错的那几句做成卡片，练三天。
\end{ansbox}

\subsection{课后野路子}

\begin{itemize}
  \item \textbf{每天饭前对口令 2 分钟}：豆角爸爸一句小豆角一句，跳着来。
  \item \textbf{生活里的口诀}：一家人吃饭几双筷子（2 $\times$ 人数）、
        一箱牛奶几盒（按排数算）、楼梯一层几级 $\times$ 几层。
  \item \textbf{把总错的 3--5 句写在卡片上贴镜子}，刷牙时念一遍。
        \textbf{只贴错的，不贴整张表。}
  \item \textbf{买东西前算总价}：一包 4 元，买 5 包多少钱？让小豆角先报口诀再报钱数。
\end{itemize}
